\section{The Minimal $1$-NN Separation}
\label{sec:minimal-1nn}

This section answers: how small can the LOO/replace-one separation be?  The
answer: two points suffice.

\subsection{Two-point construction}

Let the metric space contain two vertices $a,b$ joined by one edge, so
$d(a,b)=1$.  Consider the ordered sample
\[
S = \big((a,0),\;(b,0)\big).
\]

\subsection{LOO stability}

Delete the first occurrence $(a,0)$.  The remaining sample is $((b,0))$, whose
$1$-NN prediction at the deleted occurrence's location $a$ is still $0$ because
the sole remaining neighbor has label $0$.  Hence the loss at $(a,0)$ is
unchanged: $\Delta_{\mathrm{loo}}(S,1)=0$.

Delete the second occurrence $(b,0)$.  The remaining sample is $((a,0))$, whose
$1$-NN prediction at $b$ is $0$, unchanged.  Hence
$\Delta_{\mathrm{loo}}(S,2)=0$.

Thus $\Delta_{\mathrm{loo}}(S,i)=0$ for every sample index $i$.

\subsection{Replace-one instability}

Now evaluate at the fixed query-label pair $(x,y)=(a,0)$.  Replace the first
occurrence $(a,0)$ by $(a,1)$, obtaining
\[
S^{1 \leftarrow (a,1)} = \big((a,1),\;(b,0)\big).
\]

At query $a$, the original sample predicts $0$ because $(a,0)$ is the unique
nearest occurrence.  In the perturbed sample, the unique nearest occurrence is
now $(a,1)$, so the prediction becomes $1$.  Hence
\[
\Delta_{\mathrm{rep}}(S, 1, (a,1), a, 0) = 1.
\]

\subsection{Moral}

The separation does not require large samples, high dimension, or probabilistic
effects.  It appears already in a two-point finite metric space.  The
mechanism is minimal: a single replacement flips the label of the decisive
nearest neighbor, while each LOO deletion removes a point whose label was
already consistent with the remaining nearest neighbor.

This two-point example sets the pattern for the general constructions that
follow.  For larger $k$, the same principle (stabilize the margin on one side
for LOO, cross it via a targeted replacement) is amplified through duplicate
occurrence patterns.
